\subsection{Relative class-two approximation at coefficients at least five}

\begin{result}{E.152}{Proved}{Detection of the dual triangular module by nonexceptional Frobenius classes}
Let \(\ell\ge5\), \(e\ge1\), \(R=\mathbb Z/\ell^e\), and
let \(b\in Z^1(G_{\mathbb Q},R(1))\) be a Kummer cocycle
whose reduction modulo \(\ell\) is nonzero. Write \(\chi\)
for the cyclotomic character modulo \(\ell^e\). Define
\[
 V_b(g)=\begin{pmatrix}\chi(g)&0\\-\chi(g)b(g)&\chi(g)^2\end{pmatrix},
 \qquad
 V_b^D(g)=\begin{pmatrix}1&b(g)/\chi(g)\\0&\chi(g)^{-1}\end{pmatrix},
 \tag{E.152.1}
\]
where \(V_b^D=\operatorname{Hom}(V_b,\mu_{\ell^e})\).
The intersection of the kernels of the local restriction maps
\[
 H^1(G_{\mathbb Q,S},V_b^D)\longrightarrow
 H^1(\mathbb Q_q,V_b^D),\qquad
 q\notin S,\quad q\not\equiv\pm1\pmod\ell,
 \tag{E.152.2}
\]
is zero, for every finite \(S\) containing \(\ell,\infty\)
and all ramification of \(b\).
A finite subset of these primes already gives an injective map,
and that subset may avoid any prescribed finite set.

There is a more precise version retaining auxiliary arithmetic data.
Choose a rational representative \(B\) of \(b\), put
\(K=\mathbb Q(\mu_{\ell^e})\) and
\(K_b=K(B^{1/\ell^e})\), and let \(D/\mathbb Q\) be a finite
Galois extension such that
\[
 K\subset D,\qquad D\cap K_b=K.
 \tag{E.152.3}
\]
Include its ramification in \(S\). For
\(\sigma\in\operatorname{Gal}(D/\mathbb Q)\) with
\(\chi(\sigma)\not\equiv\pm1\pmod\ell\), restrict (E.152.2)
to primes with Frobenius conjugacy class \([\sigma]\) in \(D\).
Its kernel is exactly the inflated subgroup
\[
 \{\psi\in\operatorname{Hom}(\operatorname{Gal}(D/\mathbb Q),R):
                \psi(\sigma)=0\}\,e_1
       \subset H^1(G_{\mathbb Q,S},V_b^D).
 \tag{E.152.4}
\]
Again a finite subset of the permitted primes has this exact kernel.
\end{result}
\begin{proof}
The matrices in (E.152.1) multiply correctly because
\(b(gh)=b(g)+\chi(g)b(h)\). Their displayed duality is
\(V_b^D(g)=\chi(g)V_b(g)^{-t}\).
Since \(b\bmod\ell\ne0\), the rational number \(B\) has a
valuation prime to \(\ell\). Its real \(\ell\)-th root cannot
lie in an abelian extension of \(\mathbb Q\): its degree-\(\ell\)
real field is nonnormal. Consequently
\([K_b:K]=\ell^e\), and the image of \((b,\chi)\) is the full
semidirect product \(R\rtimes R^\times\).
In particular the translation coordinate \(b\) can be prescribed
arbitrarily above each cyclotomic value.

Write a cocycle in \(V_b^D\) as \(z(g)=(r(g),s(g))\), and put
\(\lambda=\chi(g)\), \(u=b(g)\). If
\(\lambda\not\equiv1\pmod\ell\), its class on the procyclic
group generated by \(g\) is zero precisely when
\[
 L_gz(g)=0,
 \qquad L_g=(1-\lambda,-u).
 \tag{E.152.5}
\]
Indeed the image of \(V_b^D(g)-1\) is the rank-one direct
summand generated by \((u,1-\lambda)\), up to the invertible
factor \(\lambda^{-1}\). Direct multiplication also gives
\(L_gV_b^D(g)=L_g\).

We prove the version with \(D\) first. Let
\(N=\operatorname{Gal}(\overline{\mathbb Q}/DK_b)\).
Suppose that all the permitted local restrictions of a class vanish.
Chebotarev, applied to a finite field through which its cocycle and
action factor, implies (E.152.5) for every \(g\) whose image in
\(D\) is \(\sigma\). This condition is invariant under changing
\(g\) by conjugacy, as is restriction of a cohomology class to
its cyclic subgroup. For \(n\in N\), the elements \(g\) and
\(gn\) have the same arithmetic and translation coordinates, so
\[
 0=L_g\bigl(z(gn)-z(g)\bigr)=L_gz(n).
 \tag{E.152.6}
\]
By (E.152.3), above \(\sigma\) we may take \(u=0\) and
\(u=1\). Since \(1-\lambda\) is a unit, these two rows force
\(z(n)=0\). The cocycle therefore descends to
\[
 J=\operatorname{Gal}(DK_b/\mathbb Q)
          \simeq U\rtimes\operatorname{Gal}(D/\mathbb Q),
 \qquad U\simeq R.
 \tag{E.152.7}
\]
The semidirect splitting here is the fibre-product splitting
\(d\mapsto((0,\chi(d)),d)\), not an additional assertion about
an arbitrary extension.

On \(V_b^D\), a generator of \(U\) is the matrix
\(t=\left(\begin{smallmatrix}1&1\\0&1\end{smallmatrix}\right)\).
Its norm is zero: \(\ell^e I+\binom{\ell^e}{2}(t-I)=0\) in
\(\operatorname{Mat}_2(R)\). Thus
\[
 (V_b^D)^U=Re_1,
 \qquad H^1(U,V_b^D)=V_b^D/(t-1)V_b^D\simeq Re_2.
 \tag{E.152.8}
\]
The action of an arithmetic element of cyclotomic value \(\lambda\)
on the second group is multiplication by \(\lambda^{-2}\):
one factor comes from conjugating the translation generator and
the other from its action on \(e_2\). The chosen \(\sigma\)
has \(\lambda^2-1\) a unit. Hence this cohomology group has no
arithmetic invariants. Inflation--restriction now gives
\[
 H^1(J,V_b^D)=H^1(\operatorname{Gal}(D/\mathbb Q),Re_1)
            =\operatorname{Hom}(\operatorname{Gal}(D/\mathbb Q),R)e_1.
 \tag{E.152.9}
\]
For such a class, (E.152.5) is exactly
\((1-\lambda)\psi(\sigma)=0\), hence \(\psi(\sigma)=0\).
This proves both inclusions in (E.152.4).

For (E.152.2), the same argument with the kernel of \((b,\chi)\)
first descends a class invisible at every nonexceptional Frobenius
to \(R\rtimes R^\times\). The preceding calculation identifies
its first cohomology with \(\operatorname{Hom}(R^\times,R)e_1\).
Such a character vanishes on the Teichm\"uller subgroup, whose
order is prime to \(\ell\). Fix a Teichm\"uller unit
\(\lambda_0\not\equiv\pm1\pmod\ell\). For every
\(v\in1+\ell R\), the value \(\lambda_0v\) is still
nonexceptional. Equations (E.152.5) at these values give
\(\psi(v)=\psi(\lambda_0v)=0\). Thus \(\psi=0\).
Finally \(H^1(G_{\mathbb Q,S},V_b^D)\) is finite; successively
choose a detecting Frobenius for each remaining nonzero coset of
the stated kernel. Chebotarev supplies infinitely many primes for
each choice, so finitely many suffice and any fixed finite set can
be avoided.
\end{proof}

\begin{result}{E.153}{Proved}{Finite relative class-two approximation using only ordinary new inertia}
Retain E.152. Given arbitrary local classes in
\(H^1(\mathbb Q_v,V_b)\) for \(v\in S\), they are the
restrictions of a global class in
\(H^1(G_{\mathbb Q,S\cup T},V_b)\) for some finite set
\(T\) of arbitrarily large primes
\(q\not\equiv\pm1\pmod\ell\).

Under (E.152.3), require in addition that every prime of \(T\)
have Frobenius class \([\sigma]\) in \(D\). This remains
possible whenever the first coordinates of the prescribed local
classes are the localizations of one global class
\[
                       a_0\in H^1(G_{\mathbb Q,S},R(1)).
 \tag{E.153.1}
\]
More generally the exact obstruction is the sum of local Tate
pairings with the characters in (E.152.4).

In symmetric Heisenberg coordinates, these global classes are
actual finite class-two sections with the fixed second Kummer
cocycle \(b\). At each newly ramified prime \(q\in T\), their
local section class is the image of an ordinary point
\(t_q=q^{m_q}\), for some positive integer \(m_q\).
Thus no new coefficient-\(\ell\) support-matching condition is
introduced at these primes. This is a finite-coefficient relative
approximation result; it does not by itself produce integral
class-two lifts, a compatible coefficient sequence, or an ordinary
completed pair.
\end{result}
\begin{proof}
Let \(H(R)\) have multiplication
\[
 (a,b,z)(a',b',z')=
 (a+a',b+b',z+z'+(ab'-ba')/2),
 \tag{E.153.2}
\]
with Galois weights \(1,1,2\). The actual tangential action on
this class-two quotient is this diagonal action; the conjugating
Ihara term first changes degree three. With \(b\) fixed, put
\(h=z-ab/2\). The section-cocycle equations become exactly
\[
 a(gg')=a(g)+\chi(g)a(g'),\qquad
 h(gg')=h(g)+\chi(g)^2h(g')-\chi(g)b(g)a(g').
 \tag{E.153.3}
\]
These are the cocycle equations for \(V_b\).
Local geometric changes of frame give the usual local classes,
so one can first align the second coordinate with the fixed
cocycle \(b\) and then apply the linear statement.

We prove the linear statement by finite Poitou--Tate duality
\citep[I.4.9--4.10, pp.~56--57]{CohMilneADT}.
Choose \(T\) so that restriction of the finite group
\(H^1(G_{\mathbb Q,S},V_b^D)\) to \(T\) has exactly the kernel
specified by E.152. A dual class over \(G_{\mathbb Q,S\cup T}\)
which is zero locally at all primes in \(T\) is unramified there
and therefore descends to \(G_{\mathbb Q,S}\).
Consequently this same kernel is the entire dual obstruction
to realizing prescribed local data at \(S\) while leaving the
local data at \(T\) free. To see this directly, use the perfect
sum of local pairings on \(S\cup T\): the annihilator of the
possible global localizations is the image of global dual
cohomology. Projecting to the \(S\)-coordinates retains only
dual classes with zero \(T\)-coordinates. This is the preceding
kernel.

For unrestricted nonexceptional primes the kernel is zero.
For the prescribed arithmetic class, a character in (E.152.4)
pairs only with the first \(R(1)\)-coordinate of \(V_b\).
Under (E.153.1), the obstruction is therefore
\[
       \sum_{v\in S}\operatorname{inv}_v(a_0\cup\psi)=0.
 \tag{E.153.4}
\]
This is global reciprocity: both classes are unramified outside
\(S\), and at those places their cup product is unramified and
has zero second cohomology. This proves existence. The same
argument without (E.153.1) gives the stated exact pairing
criterion.

At a prime \(q\in T\), the cocycle \(b\) is unramified and
\(q-1\) is a unit in \(R\), so its local class is zero.
After a local geometric change of frame, \(V_b\) is
\(R(1)\oplus R(2)\). Both unramified first-cohomology groups
vanish because \(q-1\) and \(q^2-1\) are units in \(R\).
The tame inertia contribution is
\[
 H^1(I_q,R(1))^{F_q}=R,
 \qquad H^1(I_q,R(2))^{F_q}=0.
 \tag{E.153.5}
\]
Wild inertia contributes nothing. Hence the whole local class
is uniquely specified by its first-coordinate inertia value
\(\bar m_q\in R\). Choose a positive integer lift \(m_q\),
using \(\ell^e\) when \(\bar m_q=0\). The ordinary point
\(q^{m_q}\) has precisely this local class: its complementary
coordinate lies in \(1+q\mathbb Z_q\), and the boundary model
has the indicated peripheral inertia and no remaining unramified
class. This proves the last assertion, including Frobenius
matching rather than inertia alone.
\end{proof}

\begin{result}{E.154}{Proved}{Arithmetic compatibility criterion and removal of the special coefficient-five collision}
In the prescribed-arithmetic version of E.153, condition
\(D\cap K_b=K\) is automatic whenever \(D\) contains \(K\)
and the commutator subgroup of \(\operatorname{Gal}(D/\mathbb Q)\)
has order prime to \(\ell\). Thus one can retain finite
arithmetic data with geometric groups at coefficient primes
other than \(\ell\), together with cyclotomic congruences,
provided an arithmetic element \(\sigma\) satisfying all those
data and \(\chi_\ell(\sigma)\not\equiv\pm1\pmod\ell\)
actually exists. The assertion is a compatibility criterion;
it does not discard this last existence requirement.

For the reference pair \((5/16,11/16)\), preserving its
quadratic first coordinate at a fresh \(B\)-support
prime forces \((5/q)=1\), and hence
\(q\equiv\pm1\pmod5\). Thus this particular reference pair
cannot use the nonexceptional coefficient-five auxiliary primes
of E.153 under that preservation condition.

The rational base
\[
                  t=23/64,\qquad 1-t=41/64
 \tag{E.154.1}
\]
removes this specific collision. Its support primes
\(2,23,41\) are all minus one modulo three; inertia at the
last two primes gives a Frattini basis at coefficient three.
The ordinary recursion of E.151 applies, replacing its factorial
offset eleven by forty-one, with coefficient forms
\[
                 a=23U^8,
 \qquad b=41V^8,
 \qquad c=64W^8.
 \tag{E.154.2}
\]
Its complementary squareclasses are \(41,23,-943\).
These quadratic characters impose no restriction on
\(q\bmod5\). At coefficients \(23\) and \(41\), a required
square residue class can still be chosen different from
\(\pm1\). Therefore the identified coefficient-five
quadratic obstruction is not intrinsic to the ordinary-pair
problem. Neither this replacement nor E.153 establishes a
simultaneous all-coefficient class-two matched construction.
\end{result}
\begin{proof}
The affine group \(\operatorname{Gal}(K_b/\mathbb Q)\)
has commutator subgroup exactly the translation group
\(U\simeq R\): conjugating a translation by a scalar
\(\lambda\) multiplies it by \(\lambda\), and some
\(\lambda-1\) is a unit. Any common quotient with a group
whose commutator has order prime to \(\ell\) must therefore
kill \(U\). The intersection is contained in the maximal
abelian subfield \(K\); since \(K\subset D\), it is exactly
\(K\). A compositum of finite groups with commutator orders
prime to \(\ell\) has the same property: its Galois group
embeds in their product, so its commutator embeds in the product
of their commutators. Adding an abelian cyclotomic field causes
no difficulty. In particular the criterion applies whenever
all the retained geometric and noncyclotomic arithmetic kernels
are groups of prime-power order at primes different from
\(\ell\). It does not apply merely because a field was
constructed at a different stage; its actual group is relevant.

For the coefficient-five assertion, the first Kummer coordinate
of the reference pair modulo squares is \(5\). At a new
\(B\)-point it is the class of a unit congruent to one.
Preserving the reference quadratic class therefore requires \((5/q)=1\).
Quadratic reciprocity gives \((5/q)=(q/5)\), and the nonzero
squares modulo five are exactly \(1,4\). The resulting
arithmetic Frobenius condition excludes every nonexceptional
residue at this coefficient. At an \(A\)-point its first unit
part is free, so this argument does not impose the same condition
there; in particular it does not forbid an \(A\)-only repair.

For (E.154.1), the same calculation instead involves
\(\sqrt{23}\) and \(\sqrt{41}\), whose conductors are
\(92\) and \(41\). The additional complementary squareclass
\(-943=-23\cdot41\) has conductor \(943\).
All these conductors are prime to five, so CRT permits either
nonexceptional residue modulo five while retaining the required
quadratic characters. The square subgroups modulo twenty-three
and forty-one have orders eleven and twenty, respectively;
after removing \(1\) and \(-1\), each still contains elements.
Thus their own quadratic restrictions do not reproduce the
coefficient-five collision.

We spell out the changes to E.151 to verify that the new base
is legitimate. Replace the private primes five and eleven by
twenty-three and forty-one, and the denominator sixteen by
sixty-four. Use \((n+41)!\) in place of \((n+11)!\), so
the radical exponent always contains the new quadratic conductors.
At three the two base units are both minus one
modulo three, so the specific-point root argument and the
choices of three, five or seven factors are unchanged.
At two there are still two odd coefficients; Newton may use
one of them. At each private odd prime there is again a unit
coefficient. The new radical ratios are
\(64W^8/(41V^8)\), \(64W^8/(23U^8)\), and
\(-41V^8/(23U^8)\). Their private valuations are \(1\) or
\(-1\), so no odd-prime radical can lie in an abelian field.
After removing rational squares and cyclotomic factors, a
fourth root would put respectively
\(41^{1/4}\), \(23^{1/4}\), or \(943^{1/4}\) in an
abelian field, again impossible because the corresponding
real quartic fields are nonnormal. Hence the radical degree
is still \(h/2\). Enlarge the fixed cyclotomic modulus to
contain the displayed quadratic conductors. The cyclic
three-primary intersection and fixed-root argument are exactly
those proved in E.151. Inertia at twenty-three and forty-one
has independent abelian coordinates, providing its required
finite geometric image. Finally the global plus line may be
pinned at two, since \((2^4-1)/3=2\bmod3\).
All remaining local and convergence arguments use only the
listed properties. This proves the claimed replacement of
the conditional ordinary recursion, with its infinite gate
still open.
The finite certificate \path{code/ghost_audit_relative_approximation.py}
checks the displayed matrix identities at coefficients five, seven,
twenty-five and forty-nine, the exceptional coefficient-three cocycle,
and the stated quadratic and CRT conditions. Its actual output is
\path{code/ghost_audit_relative_approximation.out}. These computations
check the algebra and scope; the approximation and arithmetic existence
statements use the separate proofs above.
\end{proof}
